\documentclass{article}

% Language setting
% Replace `english' with e.g. `spanish' to change the document language
%\usepackage[english,portuguese]{babel}
\usepackage[brazil]{babel}
\usepackage[utf8]{inputenc}

% Set page size and margins
% Replace `letterpaper' with `a4paper' for UK/EU standard size
\usepackage[a4paper,top=2cm,bottom=2cm,left=3cm,right=3cm,marginparwidth=1.75cm]{geometry}

% Useful packages
\usepackage{amsmath,amsthm,amsfonts,float}
\usepackage{graphicx}
\usepackage[colorlinks=true, allcolors=blue]{hyperref}
\usepackage{multicol}

%\newenvironment{namebox}
%{\begin{flushleft}Nome: \underline{\hspace{6cm}}\end{flushleft}} % Adjust the length of the underline as needed

\newcommand{\na}{\mathbb{N}} % Set of integers
\newcommand{\naz}{\mathbb{N}_0} % Set of integers
\newcommand{\Z}{\mathbb{Z}} % Set of integers
\newcommand{\re}{\mathbb{R}} % Set of real numbers
\newcommand{\sinal}{\mathrm{sinal}} % Signal operador

\newcommand{\vp}{{\mathbf p}}
\newcommand{\vP}{{\mathbf P}}

\newcommand{\vpi}{{\boldsymbol \pi}}
%\newcommand{\Pr}{{\mathbb P}}


\date{--/--/----} % Removendo a data

\begin{document}

\begin{flushleft}
	Universidade Federal da Bahia\\
	Instituto de Matemática e Estatística\\
	Departamento de Estatística\\
	MATF14 - Estatística Econômica I - 2026.1\\
	Professor: Rodney Fonseca \\
%	Data: 24/09/2024
\end{flushleft}

\begin{center}
	\begin{Large}
		Lista 4 de exercícios \\
	\end{Large}	
%	Prazo de entrega: 4/11/2024
\end{center}

%\fbox{\begin{minipage}{30em}
%		Nome:
%\end{minipage}}
%\hspace{2em} Nota: 

%\vspace{2em}

\vspace{1em}

\section*{Conceitos elementares de cálculo}

\paragraph{Derivada da função constante:} $f(x) = c \Rightarrow f'(x) = 0$

\paragraph{Derivada da função potência:} $f(x) = x^n \Rightarrow f'(x) = n x^{n-1}$

\paragraph{Derivada da função exponencial:} $f(x) = e^x \Rightarrow f'(x) = e^x$

\paragraph{Derivada da função logarítmica:} para $x>0$, $f(x) = \ln x \Rightarrow f'(x) = 1/x$

\paragraph{Derivada da soma:} $f(x) = u(x) + v(x) \Rightarrow f'(x) = u'(x) + v'(x)$

\paragraph{Derivada do produto:} $f(x) = u(x) v(x) \Rightarrow f'(x) = u'(x)v(x) + u(x) v'(x)$

\paragraph{Derivada da função vezes uma constante:} $f(x) = c \cdot u(x) \Rightarrow f'(x) = c \cdot u'(x)$

\paragraph{Regra da cadeia:} $f(x) = g(u(x)) \Rightarrow f'(x) = g'(u(x)) \cdot u'(x)$

\paragraph{Cálculo de integrais:} $F$ é a {\em primitiva} de $f$ se $F'(x) = f(x)$. Assim, dado um intervalo $(a,b)$, 
$$
\int_a^b f(x)dx = F(b) - F(a)
$$


\vspace{1em}

\section*{Exercícios}

\begin{enumerate}

    
\item Obtenha a derivada das seguintes funções:\\
$f(x) = 5, \quad g(x)=x^6, \quad h(x)=x^{15}$
   
\item Calcule a derivada de cada uma das seguintes funções:
\begin{enumerate}
\begin{multicols}{2}
\item $f(x) = 8 x^{11}$
\item $f(x) = -\frac{7}{5}x^3 - \frac{\sqrt{3}}{7}$
\item $f(x) = 5 + x + 3x^2$
\item $f(x) = 3 + 2x^n + x^{2n}$, com $n\in\mathbb{N}$ fixado
\item $f(x) = e^{-3x}$
\item $f(x) = e^{7x^2 - 2x}$
\end{multicols}
\end{enumerate}

   
\item Determine a primitiva das seguintes funções:
\begin{enumerate}
\begin{multicols}{3}
\item $f(x) = \sqrt{x}$
\item $f(x) = \frac{1}{x^3}$
\item $f(x) = x^{-2/5}$
\end{multicols}
\end{enumerate}

   
\item Calcule as integrais definidas abaixo:
\begin{enumerate}
\begin{multicols}{3}
\item $\int_0^1 7 dx$
\item $\int_1^2 x^3 dx$
\item $\int_0^1 (-x^2) dx$
\item $\int_0^2 (x^2 - 3x + 5) dx$
\item $\int_0^2 x(x^5 - 1) dx$
\item $\int_0^3 e^{2x} dx$
\end{multicols}
\end{enumerate}

	
		
	
\end{enumerate}


\end{document}